\documentclass[12pt,a4paper]{article}

% -------------------------------------------------
% Pacotes
% -------------------------------------------------
\usepackage[utf8]{inputenc}
\usepackage[brazil]{babel}
\usepackage{geometry}
\usepackage{graphicx}
\usepackage{hyperref}
\usepackage{float}
\usepackage{setspace}

\geometry{top=2.5cm,bottom=2.5cm,left=2.5cm,right=2.5cm}
\onehalfspacing

% -------------------------------------------------
% Título
% -------------------------------------------------
\title{
Apontamento de Literatura – Tema 1\\
\large Metodologia de Investigação em Revisão Sistemática
}
\author{
Disciplina: Revisão Sistemática e Meta-análises\\
Pós-graduação (Mestrado e Doutoramento)
}
\date{}

\begin{document}

\maketitle

% =================================================
\section{Introdução à Investigação Científica}
% =================================================

A investigação científica constitui um processo sistemático, rigoroso e cumulativo de produção de conhecimento, orientado pela observação empírica, formulação teórica e validação metodológica. No contexto da revisão sistemática, a investigação assume um papel central na organização, síntese e avaliação crítica da evidência existente.

Os principais objetivos da investigação científica incluem:
\begin{enumerate}
    \item Descrever fenómenos naturais ou sociais;
    \item Explicar relações causais ou associativas entre variáveis;
    \item Prever comportamentos ou resultados futuros;
    \item Sustentar a tomada de decisão baseada na evidência.
\end{enumerate}

A pesquisa científica distingue-se de abordagens narrativas ou opinativas por sua reprodutibilidade, transparência metodológica e fundamentação empírica.

% =================================================
\section{Pirâmide dos Níveis de Evidência}
% =================================================

A hierarquização dos níveis de evidência é fundamental para a avaliação da robustez científica dos resultados de investigação. Diferentes desenhos de estudo produzem evidência com distintos graus de confiabilidade, suscetibilidade a viés e capacidade de inferência causal.

De forma geral, a pirâmide de evidência organiza-se da seguinte maneira (da base ao topo):
\begin{itemize}
    \item Opinião de especialistas;
    \item Relatos e séries de casos;
    \item Estudos transversais e observacionais;
    \item Estudos de coorte e caso-controle;
    \item Ensaios clínicos randomizados;
    \item Revisões sistemáticas e meta-análises.
\end{itemize}

\begin{figure}[H]
    \centering
    \includegraphics[width=0.65\textwidth]{piramide_evidencia.png}
    \caption{Pirâmide dos níveis de evidência em investigação científica}
    \label{fig:piramide}
\end{figure}

As revisões sistemáticas e meta-análises ocupam o nível mais elevado, pois sintetizam resultados de múltiplos estudos primários mediante critérios explícitos e reprodutíveis.

% =================================================
\section{Tipos de Desenhos de Investigação}
% =================================================

Os desenhos de investigação podem ser classificados em quantitativos, qualitativos e mistos, de acordo com a natureza dos dados e os objetivos analíticos.

\begin{table}[H]
\centering
\begin{tabular}{|p{3cm}|p{6cm}|p{4cm}|}
\hline
\textbf{Tipo} & \textbf{Características} & \textbf{Exemplos} \\
\hline
Quantitativo & Dados numéricos, hipóteses testáveis, inferência estatística & Ensaios clínicos, estudos de coorte \\
\hline
Qualitativo & Exploração interpretativa de fenómenos complexos & Entrevistas, grupos focais \\
\hline
Misto & Integração de métodos quantitativos e qualitativos & Surveys com entrevistas \\
\hline
\end{tabular}
\caption{Classificação dos desenhos de investigação}
\end{table}

No contexto de revisões sistemáticas, estes desenhos determinam critérios de inclusão, métodos de síntese e tipos de análise.

% =================================================
\section{Elementos Fundamentais de um Estudo de Investigação}
% =================================================

Independentemente do desenho metodológico, um estudo científico robusto deve contemplar os seguintes elementos fundamentais:
\begin{enumerate}
    \item Formulação clara do problema de pesquisa;
    \item Definição explícita de objetivos e/ou hipóteses;
    \item Identificação das variáveis relevantes;
    \item Descrição da população e da amostra;
    \item Metodologia transparente e reprodutível;
    \item Estratégia de análise adequada;
    \item Interpretação crítica dos resultados.
\end{enumerate}

A ausência ou fragilidade de qualquer destes elementos compromete a validade interna e externa do estudo.

% =================================================
\section{Medidas de Força em Investigação Quantitativa}
% =================================================

As medidas de força quantificam associações, diferenças ou efeitos entre variáveis. As mais utilizadas incluem:
\begin{itemize}
    \item Risco Relativo (RR);
    \item Odds Ratio (OR);
    \item Coeficiente de correlação;
    \item Tamanho do efeito (Cohen’s d, Hedges’ g).
\end{itemize}

Estas medidas constituem a base estatística das meta-análises e permitem a combinação quantitativa de resultados de estudos independentes.

% =================================================
\section{Investigação Quantitativa vs. Qualitativa}
% =================================================

A investigação quantitativa e qualitativa diferem em pressupostos epistemológicos, metodologias e formas de análise.

\begin{table}[H]
\centering
\begin{tabular}{|p{4cm}|p{4cm}|p{4cm}|}
\hline
\textbf{Aspecto} & \textbf{Quantitativa} & \textbf{Qualitativa} \\
\hline
Tipo de dados & Numéricos & Textuais ou visuais \\
\hline
Objetivo & Testar hipóteses & Explorar significados \\
\hline
Análise & Estatística & Temática / interpretativa \\
\hline
\end{tabular}
\caption{Comparação entre investigação quantitativa e qualitativa}
\end{table}

\begin{figure}[H]
    \centering
    \includegraphics[width=0.7\textwidth]{comparacao_quant_qual.png}
    \caption{Comparação conceptual entre abordagens quantitativa e qualitativa}
    \label{fig:comparacao}
\end{figure}

% =================================================
\section{Avaliação Crítica de um Artigo Científico}
% =================================================

A avaliação crítica de artigos científicos constitui uma competência central em revisões sistemáticas. Recomenda-se que a análise considere, pelo menos, as seguintes questões:

\begin{enumerate}
    \item Qual é a questão de investigação?
    \item O objetivo está claramente definido?
    \item O desenho do estudo é apropriado?
    \item A amostra é adequada e representativa?
    \item Os métodos são válidos e confiáveis?
    \item A análise é estatisticamente apropriada?
    \item Os resultados são apresentados de forma clara?
    \item As conclusões são sustentadas pelos dados?
    \item As limitações são discutidas?
    \item O estudo contribui para o avanço do conhecimento?
\end{enumerate}

Esta estrutura analítica fornece a base para decisões de inclusão, exclusão e ponderação de estudos em revisões sistemáticas.

\end{document}
